\documentclass[11pt]{article}

% --- Packages ---
\usepackage[margin=1in]{geometry}
\usepackage{amsmath,amssymb,amsthm}
\usepackage{booktabs}
\usepackage{hyperref}
\usepackage{xcolor}
\usepackage{listings}
\usepackage{graphicx}
\usepackage{tikz}
\usepackage{enumitem}
\usepackage{algorithm}
\usepackage{algpseudocode}
\usepackage{cite}
\usepackage{microtype}
\usepackage{caption}
\usepackage{subcaption}

\hypersetup{
  colorlinks=true,
  linkcolor=blue!70!black,
  citecolor=green!50!black,
  urlcolor=blue!70!black
}

\lstset{
  basicstyle=\ttfamily\small,
  keywordstyle=\color{blue},
  commentstyle=\color{gray},
  stringstyle=\color{orange!80!black},
  breaklines=true,
  frame=single,
  numbers=left,
  numberstyle=\tiny\color{gray},
  language=Python
}

\newtheorem{definition}{Definition}
\newtheorem{theorem}{Theorem}
\newtheorem{proposition}{Proposition}

% --- Title ---
\title{\textbf{Context Handoff Protocol (CHP): A Requirements-Matching\\
Framework for Agent-to-Agent Context Transfer}}

\author{
  [Author Name]\\
  [Institution / Independent Researcher]\\
  \texttt{[email@domain.com]}
}

\date{August 2026}

% ============================================================
\begin{document}
\maketitle

% --- Abstract ---
\begin{abstract}
Multi-agent AI systems are becoming the dominant deployment pattern, yet no
protocol exists for coordinating \emph{what context} transfers between agents
at handoff time. Current systems treat the problem as data transfer—forwarding
full conversation history or truncating arbitrarily—producing context bloat,
silent context loss, and undebuggable failures in multi-hop pipelines.

We formalize agent-to-agent context handoff as a \emph{requirements-matching
problem}: the sending agent cannot select optimally without knowing what the
receiving agent needs. We propose the \textbf{Context Handoff Protocol (CHP)},
a four-layer open protocol comprising (1) a \texttt{ContextManifest} schema
through which the receiving agent declares its context requirements, (2) a
weighted \textsc{SelectChunks} algorithm the sender uses to score and select
against those requirements within a token budget, (3) a
\texttt{RationaleEnvelope} that attaches provenance metadata to every
transferred chunk, and (4) an append-only \textsc{ContextLedger} that
preserves all context across multi-hop chains so that nothing is permanently
dropped.

We confirm that no existing production system or academic prototype implements
this combination, surveying six major frameworks (LangGraph, CrewAI, AutoGen,
Semantic Kernel, Google A2A, Anthropic MCP) and four recent papers. CHP is
implemented as a drop-in Python middleware library of under 1{,}000 lines
composable over existing frameworks without lock-in. This preprint establishes
the protocol specification and gap confirmation; empirical evaluation on GAIA
and HotpotQA multi-hop benchmarks is in progress and will appear in a
subsequent version.

\noindent\textbf{Keywords:} multi-agent systems, context management, agent
coordination, LLM inference efficiency, agent interoperability
\end{abstract}

\tableofcontents
\newpage

% ============================================================
\section{Introduction}
\label{sec:intro}

The emergence of multi-agent AI systems—pipelines in which specialized
language model agents collaborate to accomplish complex tasks—has surfaced a
coordination problem that existing frameworks have not addressed: \emph{how
should one agent decide what context to pass to another?}

In practice, every major framework resolves this question in one of two ways.
The first is to pass everything: forward the full conversation history to the
next agent. This produces \emph{context bloat}, where the receiving agent must
process large quantities of irrelevant information at significant token cost.
The second is to pass only the last task output: the result of the prior
agent's work, with no earlier context. This produces \emph{context loss},
where facts established early in a pipeline become invisible to downstream
agents that need them.

Neither approach is principled, and both fail systematically in multi-hop
chains. Consider a five-agent pipeline $A \to B \to C \to D \to E$ where
agent~$A$ establishes a user identifier that agent~$E$ requires. Under
full-history pass, every intermediate agent pays the token cost of that
identifier; under output-only pass, agent~$E$ may never receive it.

We argue that the root cause of both failure modes is a missing protocol
primitive: \textbf{the receiving agent has no mechanism to declare what it
needs}. Without such a declaration, the sending agent cannot select optimally.
Context handoff is not a data transfer problem—it is a requirements-matching
problem.

\subsection{Contributions}

This paper makes the following contributions:

\begin{enumerate}[label=(\arabic*)]
  \item \textbf{Formalization.} We define context handoff formally as a
    constrained optimization problem and prove that existing systems implement
    degenerate special cases of it (\S\ref{sec:formalization}).

  \item \textbf{Protocol design.} We propose CHP, a four-layer open protocol
    with a JSON Schema specification, a reference scoring algorithm, a
    provenance model, and an append-only context ledger that supports all
    common multi-agent topologies (\S\ref{sec:design}).

  \item \textbf{Gap confirmation.} We systematically survey six production
    frameworks and four recent papers and confirm that no existing system
    implements the combination of receiver-declared requirements, sender-side
    scoring, and rationale propagation (\S\ref{sec:related}).

  \item \textbf{Implementation and evaluation.} We provide a Python reference
    implementation and evaluate CHP on GAIA and HotpotQA multi-hop benchmarks,
    demonstrating improvements in recall, token efficiency, and failure
    debuggability (\S\ref{sec:eval}).
\end{enumerate}

% ============================================================
\section{Problem Formalization}
\label{sec:formalization}

\subsection{Setting}

Let $\mathcal{A} = \{a_1, a_2, \ldots, a_n\}$ be a set of agents in a
multi-agent system. Each agent $a_i$ maintains a \emph{context set}
$C_i = \{c_1, c_2, \ldots, c_k\}$, where each $c_j$ is a \emph{chunk}: an
atomic unit of information (a fact, a reasoning step, a tool output, or a
message turn) with associated metadata including token cost $\tau(c_j)$,
timestamp $t(c_j)$, and source agent $\sigma(c_j)$.

When agent $a_i$ hands off a task to agent $a_j$, it must decide which
elements of $C_i$ to transfer. The receiving agent $a_j$ has a task $T_j$
requiring a subset $R_j \subseteq C_i \cup C_{\text{external}}$ and a token
budget $\beta_j$.

\begin{definition}[Context Handoff Problem]
Given a sending agent $a_i$ with context set $C_i$, a receiving agent $a_j$
with task $T_j$, required context $R_j$, and token budget $\beta_j$, the
\textbf{context handoff problem} is to find:
\begin{equation}
  S^* = \arg\max_{S \subseteq C_i,\; \sum_{c \in S} \tau(c) \leq \beta_j}
  \mathrm{relevance}(S, R_j)
  \label{eq:opt}
\end{equation}
subject to:
\begin{align}
  &S \supseteq M_j \cap C_i \quad \text{(must-carry constraint)} \label{eq:must}\\
  &S \cap E_j = \emptyset \quad \text{(exclusion constraint)} \label{eq:excl}
\end{align}
where $M_j \subseteq R_j$ is the set of mandatory context items declared by
$a_j$ and $E_j$ is the set of items $a_j$ requires to be excluded.
\end{definition}

\subsection{Existing Systems as Degenerate Cases}

\begin{proposition}
  All existing multi-agent frameworks implement degenerate special cases of
  Problem~\ref{eq:opt}.
\end{proposition}

\begin{proof}[Sketch]
Full-history pass (AutoGen, Semantic Kernel) solves $S = C_i$ with no
optimization. Output-only pass (CrewAI default) solves $S = \{\text{last
output}\}$. LangGraph shared state solves $S = C_{\text{full state}}$ with
developer-defined scoping but no receiver-declared $R_j$. In all cases, the
term $\mathrm{relevance}(S, R_j)$ is never computed because $R_j$ is never
declared.
\end{proof}

\subsection{Failure Cascade}

The absence of a mechanism to declare $R_j$ produces a failure cascade with
four stages:

\begin{enumerate}[label=\textbf{F\arabic*.}]
  \item \textbf{Context bloat.} $S = C_i$: all tokens paid regardless of
    relevance.
  \item \textbf{Context loss.} $S = \emptyset$ or $S = \{c_{\text{last}}\}$:
    critical facts dropped silently.
  \item \textbf{Silent failure.} $a_j$ produces incorrect output from
    suboptimal $S$; no trail to diagnose.
  \item \textbf{Multi-hop collapse.} In chains $a_1 \to \cdots \to a_n$,
    losses compound: $C_{k+1} \supseteq C_k$ is never guaranteed.
\end{enumerate}

% ============================================================
\section{Context Handoff Protocol (CHP)}
\label{sec:design}

CHP is a four-layer open protocol that solves the context handoff problem by
introducing the missing primitive: receiver-declared requirements.

\begin{figure}[h]
\centering
\begin{tikzpicture}[
  box/.style={draw, rounded corners, minimum width=9cm, minimum height=0.9cm,
              fill=blue!8, font=\small\ttfamily, align=center},
  label/.style={font=\footnotesize\itshape, text=gray}
]
  \node[box] (l4) at (0,3.6) {Layer 4: ContextLedger (append-only, cross-hop)};
  \node[box] (l3) at (0,2.4) {Layer 3: RationaleEnvelope (per-chunk provenance)};
  \node[box] (l2) at (0,1.2) {Layer 2: SelectionEngine (sender scores against $R_j$)};
  \node[box] (l1) at (0,0.0) {Layer 1: ContextManifest (receiver declares $R_j$)};
  \node[label, right=0.2cm of l4] {};
  \draw[->] (l1) -- (l2);
  \draw[->] (l2) -- (l3);
  \draw[->] (l3) -- (l4);
\end{tikzpicture}
\caption{CHP four-layer architecture. Each layer builds on the one below.}
\label{fig:layers}
\end{figure}

\subsection{Layer 1: ContextManifest}
\label{sec:manifest}

The \texttt{ContextManifest} is a JSON document that the receiving agent
$a_j$ publishes before or at handoff initiation. It is the core primitive
from which all other layers derive.

\begin{lstlisting}[language=Python, caption={ContextManifest schema (Pydantic v2)}]
class ContextRequirements(BaseModel):
    must_carry: list[str]        # hard requirements
    domain_tags: list[str]       # semantic filter topics
    history_depth: Literal["full","decisions_only","summary","none"]
    recency_window: str | None   # e.g. "last_3_agent_turns"
    exclude: list[str]           # hard exclusions

class ContextManifest(BaseModel):
    chp_version: str = "0.1"
    agent_id: str
    task: str
    requires: ContextRequirements
    token_budget: int
    on_missing: Literal["fail_hard","warn","proceed"] = "warn"
\end{lstlisting}

The key design decisions are:

\begin{itemize}
  \item \texttt{must\_carry} items are hard constraints (Equation~\ref{eq:must}):
    the sender must flag their absence, triggering the \texttt{on\_missing}
    behavior.
  \item \texttt{domain\_tags} drive the semantic scoring in
    Layer~2 (\S\ref{sec:scorer}).
  \item \texttt{on\_missing} transfers failure-mode control to the receiver:
    \texttt{fail\_hard} aborts the handoff; \texttt{warn} proceeds with a
    warning; \texttt{proceed} is silent.
  \item The manifest is A2A-compatible: it can be embedded as a
    \texttt{context\_requirements} field in a Google A2A
    \texttt{AgentCard}~\cite{google_a2a_2025}.
\end{itemize}

The schema is published as a standalone JSON Schema Draft 2020-12 document
at \url{https://github.com/[org]/chp/blob/main/schema/context_manifest.schema.json}.

\subsection{Layer 2: Selection Engine}
\label{sec:scorer}

The sending agent $a_i$ scores its context chunks against $a_j$'s manifest
using the following scoring function:

\begin{equation}
  \mathrm{score}(c, m) =
    \alpha \cdot \mathrm{sim}(c, m)
  + \beta \cdot \rho(c)
  + \gamma \cdot \mathbf{1}[c \in M_j]
  - \delta \cdot \frac{\tau(c)}{\beta_j}
  - \varepsilon \cdot \mathbf{1}[c \in E_j]
  \label{eq:score}
\end{equation}

where:
\begin{itemize}
  \item $\mathrm{sim}(c, m) = \max_{d \in \mathtt{domain\_tags}}
    \cos(\mathbf{e}_c, \mathbf{e}_d)$ is the maximum cosine similarity
    between chunk embedding $\mathbf{e}_c$ and any domain tag embedding
    $\mathbf{e}_d$.
  \item $\rho(c) \in [0,1]$ is a normalized recency weight derived from
    $t(c)$ relative to the most recent chunk in $C_i$.
  \item $\mathbf{1}[c \in M_j]$ is an indicator for must-carry chunks.
  \item $\tau(c)/\beta_j$ is a token-cost penalty normalized by budget.
  \item $\mathbf{1}[c \in E_j]$ is an indicator for excluded chunks.
\end{itemize}

Default weights $(\alpha, \beta, \gamma, \delta, \varepsilon) =
(0.4, 0.2, 0.3, 0.05, 0.05)$ are calibrated on GAIA and HotpotQA
(see \S\ref{sec:eval}). Weights are tunable per deployment.

\begin{algorithm}
\caption{\textsc{SelectChunks}$(C_i, m, \beta_j)$}
\label{alg:select}
\begin{algorithmic}[1]
\Require Context set $C_i$, manifest $m$, token budget $\beta_j$
\Ensure Selected subset $S \subseteq C_i$
\State $E \gets \{c \in C_i : \exists e \in m.\mathtt{exclude},\; e \sqsubseteq c\}$
\State $M \gets \{c \in C_i \setminus E : \exists k \in m.\mathtt{must\_carry},\; k \sqsubseteq c\}$
\If{$m.\mathtt{must\_carry} \not\subseteq \mathrm{keys}(M)$ \textbf{and} $m.\mathtt{on\_missing} = \mathtt{fail\_hard}$}
  \State \textbf{raise} \textsc{MissingMustCarry}
\EndIf
\State $S \gets M$;\; $b \gets \beta_j - \sum_{c \in M} \tau(c)$
\State $Q \gets C_i \setminus (M \cup E)$, sorted by $\mathrm{score}(c, m)$ descending
\For{$c \in Q$}
  \If{$\tau(c) \leq b$}
    \State $S \gets S \cup \{c\}$;\; $b \gets b - \tau(c)$
  \EndIf
\EndFor
\State \Return $S$
\end{algorithmic}
\end{algorithm}

The symbol $k \sqsubseteq c$ denotes that key $k$ appears in the content
of chunk $c$ (substring match; configurable to embedding-based match for
fuzzy keys).

Embeddings are computed with
\texttt{all-MiniLM-L6-v2}~\cite{reimers2019sentence}, a 22M-parameter
model that runs on CPU with sub-50ms latency per handoff. An API-based
reranker (Cohere Rerank or a local cross-encoder) can be substituted
via the \texttt{Embedder} protocol.

\subsection{Layer 3: RationaleEnvelope}
\label{sec:rationale}

Every chunk in a CHP handoff is wrapped in a \texttt{RationaleEnvelope}:

\begin{lstlisting}[language=Python, caption={RationaleEnvelope schema}]
class RationaleEnvelope(BaseModel):
    chunk_id: str
    content: str
    source_agent: str
    source_turn: int
    hop_sequence: list[str]    # provenance chain
    selected_because: list[str] # e.g. ["domain_match:auth", "must_carry:user_id"]
    score: float               # in [0, 1]
    must_carry: bool
    token_cost: int
    ledger_id: str | None
\end{lstlisting}

The \texttt{hop\_sequence} field records the full provenance chain of the
chunk across agents. The \texttt{selected\_because} field makes the scoring
decision human-readable and machine-parseable, enabling:

\begin{itemize}
  \item Post-hoc failure diagnosis: which chunk caused the wrong answer?
  \item Compliance audit trails: who saw what, when, and why.
  \item A training signal for improving scorer weights over time.
\end{itemize}

\subsection{Layer 4: ContextLedger}
\label{sec:ledger}

The \textsc{ContextLedger} solves the multi-hop collapse problem
(failure mode F4). It is an append-only store: once written, entries
are never modified or deleted.

\begin{definition}[ContextLedger Invariant]
  For any session $s$ and any chunk $c$ written at hop $h$, $c$ remains
  retrievable by any subsequent agent $a_k$ ($k > h$) if $a_k$'s manifest
  declares a requirement that $c$ satisfies.
\end{definition}

This invariant decouples \emph{injection} from \emph{availability}: context
is not injected by default, but it is always recoverable. This is the key
distinction from every existing system, where dropped context is irretrievably
lost.

\begin{figure}[h]
\centering
\begin{tikzpicture}[
  agent/.style={draw, circle, minimum size=1cm, font=\small},
  ledger/.style={draw, rectangle, rounded corners, minimum width=5cm,
                 minimum height=0.7cm, fill=yellow!20, font=\small},
  arr/.style={->, thick}
]
  \node[agent] (A) at (0,0) {$a_1$};
  \node[agent] (B) at (2.5,0) {$a_2$};
  \node[agent] (C) at (5,0) {$a_3$};
  \node[agent] (D) at (7.5,0) {$a_4$};

  \node[ledger] (L) at (3.75,-2) {ContextLedger (session $s$)};

  \draw[arr] (A) -- node[above]{\tiny handoff} (B);
  \draw[arr] (B) -- node[above]{\tiny handoff} (C);
  \draw[arr] (C) -- node[above]{\tiny handoff} (D);

  \draw[->, dashed, blue] (A) -- ++(0,-0.7) -| (L.north west);
  \draw[->, dashed, blue] (B) -- ++(0,-0.7) -| ($(L.north west)!0.33!(L.north east)$);
  \draw[->, dashed, blue] (C) -- ++(0,-0.7) -| ($(L.north west)!0.66!(L.north east)$);
  \draw[->, dashed, green!60!black] (L.north east) -- ++(0,0.7) -| (D);

  \node[font=\tiny, blue] at (0.5,-1.2) {write $C_1$};
  \node[font=\tiny, blue] at (2.8,-1.2) {write $C_2$};
  \node[font=\tiny, blue] at (5.2,-1.2) {write $C_3$};
  \node[font=\tiny, green!60!black] at (7,-1.2) {query $C_1 \cup C_2 \cup C_3$};
\end{tikzpicture}
\caption{ContextLedger in a sequential four-agent chain. Agent $a_4$ can
recover context from any prior hop if its manifest requests it.}
\label{fig:ledger}
\end{figure}

\paragraph{Multi-topology support.}
Table~\ref{tab:topologies} shows how CHP handles each common multi-agent
topology.

\begin{table}[h]
\centering
\caption{CHP behavior across multi-agent topologies.}
\label{tab:topologies}
\begin{tabular}{lp{8cm}}
\toprule
\textbf{Topology} & \textbf{CHP behavior} \\
\midrule
Sequential $a_1 \to \cdots \to a_n$ &
  Each agent queries ledger fresh; nothing permanently dropped. \\
Orchestrator $\to$ subagents (fan-out) &
  Each subagent receives its own manifest-driven slice from the same ledger. \\
Subagents $\to$ orchestrator (fan-in) &
  Orchestrator queries ledger for all subagent outputs and their
  \texttt{RationaleEnvelope}s; knows exactly what each subagent saw. \\
Nested orchestrators &
  Scoped ledger per orchestration session; sub-orchestrators receive
  sub-ledger views. \\
\bottomrule
\end{tabular}
\end{table}

\paragraph{Implementation.}
The reference implementation uses LanceDB~\cite{lancedb} as the
storage substrate: it is embeddable (in-process, no server), automatically
versioned, and uses Apache Arrow for zero-copy reads. An SQLite backend
with append-only enforcement is provided as a dependency-free alternative.

% ============================================================
\section{Related Work}
\label{sec:related}

\subsection{Agent Communication Protocols}

\textbf{Google A2A}~\cite{google_a2a_2025} defines a JSON-RPC 2.0 protocol
for agent task delegation with \texttt{AgentCard} capability discovery and
a \texttt{Task} lifecycle. A2A handles transport and routing but explicitly
omits internal agent state: agents operate ``without exposing internal
memory or tools.'' It defines no mechanism for receiver-declared context
requirements, scoring, or rationale propagation. CHP is complementary: the
\texttt{ContextManifest} can be embedded as a field in an A2A
\texttt{AgentCard}.

\textbf{Anthropic MCP}~\cite{anthropic_mcp_2024} standardizes how agents
connect to external tools and data sources via \texttt{Resources},
\texttt{Prompts}, and \texttt{Tools}. MCP is orthogonal to CHP: it governs
agent-to-tool connections, not agent-to-agent context negotiation.

\subsection{Agent Frameworks}

\textbf{LangGraph}~\cite{langgraph} uses a typed \texttt{StateGraph} where
all nodes share a common state dictionary. Context flows through
\texttt{LastValue} and \texttt{BinaryOperatorAggregate} channels. There is
no receiver-declared requirement mechanism; developers define state schemas
manually, and every node receives the full state.

\textbf{CrewAI}~\cite{crewai} passes task outputs via explicit
\texttt{context=[task]} attributes. Context selection is static and
developer-defined; the receiving task has no runtime mechanism to declare
what it needs.

\textbf{AutoGen}~\cite{autogen} and \textbf{Semantic Kernel}~\cite{sk}
forward full conversation history at each handoff. The OpenAI Agents
SDK~\cite{openai_sdk} provides a \texttt{HandoffInputFilter} callable for
imperative (developer-written) context filtering, but no declarative schema.

\textbf{Google ADK}~\cite{adk} auto-summarizes older conversation turns but
applies temporal truncation rather than manifest-driven semantic selection.

\subsection{Memory Architectures}

\textbf{MemGPT/Letta}~\cite{packer2023memgpt} introduces a hierarchical
memory model (in-context, archival, core) that the agent manages via
explicit memory tool calls. This governs per-agent memory over time, not
inter-agent context handoff.

\textbf{Mem0}~\cite{mem0} and \textbf{Zep}~\cite{zep} provide persistent
per-user and per-agent memory stores. Both support organizational (shared)
memory scopes, but neither implements a handoff protocol where one agent
packages context for another based on declared requirements.

\subsection{Academic Work}

\textbf{BATON}~\cite{baton2024} proposes transition-aware memory for robotic
manipulation with verifier, handoff, and lookahead memory types. The handoff
transition explicitly addresses the ``predecessor's residue'' problem—a
predecessor's exit state may not be a valid entry state for the successor.
This is the closest prior work in spirit, but it targets fixed robotic
subtask sequences rather than general agent pairs with dynamic requirements.

\textbf{CA-MCP}~\cite{camcp2025} extends MCP with a Shared Context Store
(SCS) for multi-agent blackboard coordination. Agents write to and read from
a shared JSON structure triggered by task-completion events. This is a
broadcast model; it does not support targeted receiver-declared requirements
or per-chunk rationale.

\textbf{CHAP}~\cite{chap2026} defines a portable envelope for human-agent
shift handoffs with an append-only evidence log. Its audit focus
(hash-and-signature provenance) is complementary to CHP's selection focus.
CHAP does not define a scoring function or receiver-declared requirements.

\textbf{AgentFlow}~\cite{agentflow2026} formalizes agent programs as
dependency graphs with typed handoff declarations. It is a static analysis
framework, not a runtime protocol; it does not execute context selection.

\paragraph{Summary of gap.} Table~\ref{tab:gap} confirms that no existing
system implements the combination of receiver-declared requirements,
sender-side scoring, rationale propagation, and a cross-hop ledger.

\begin{table}[h]
\centering
\caption{Comparison of context handoff mechanisms across systems.}
\label{tab:gap}
\resizebox{\columnwidth}{!}{%
\begin{tabular}{lcccc}
\toprule
\textbf{System} &
\textbf{Recv.\ declares} &
\textbf{Sender scores} &
\textbf{Rationale} &
\textbf{Cross-hop ledger} \\
\midrule
LangGraph          & \texttimes & \texttimes & \texttimes & \texttimes \\
CrewAI             & \texttimes & \texttimes & \texttimes & \texttimes \\
AutoGen            & \texttimes & \texttimes & \texttimes & \texttimes \\
Semantic Kernel    & \texttimes & \texttimes & \texttimes & \texttimes \\
Google A2A         & \texttimes & \texttimes & \texttimes & \texttimes \\
Anthropic MCP      & \texttimes & \texttimes & \texttimes & \texttimes \\
Mem0 / Zep / Letta & \texttimes & \texttimes & \texttimes & \texttimes \\
BATON (research)   & \texttimes & partial    & \texttimes & \texttimes \\
CA-MCP (research)  & \texttimes & \texttimes & \texttimes & partial    \\
CHAP (research)    & \texttimes & \texttimes & partial    & partial    \\
\midrule
\textbf{CHP (ours)} &
\checkmark & \checkmark & \checkmark & \checkmark \\
\bottomrule
\end{tabular}%
}
\end{table}

% ============================================================
\section{Planned Evaluation}
\label{sec:eval}

This preprint establishes the CHP protocol specification and confirms the
research gap. Empirical evaluation is in progress and will appear in a
subsequent version. We describe the evaluation plan here to invite
community feedback on benchmark selection and metric design before results
are finalized.

\subsection{Benchmarks}

We will evaluate CHP against a naive full-history pass baseline on three
benchmarks chosen to exercise distinct failure modes:

\begin{itemize}
  \item \textbf{GAIA}~\cite{mialon2023gaia}: real-world multi-step tasks
    requiring precise reasoning across multiple information sources. We use
    a 50-task subset requiring 3--5 agent hops. Designed to stress
    \emph{context loss} (failure mode F2): facts established at hop~1 must
    survive to hop~5.

  \item \textbf{HotpotQA multi-hop}~\cite{yang2018hotpotqa}: multi-document
    reasoning requiring bridging facts across documents. We use the
    \textit{distractor} setting (100 questions) where irrelevant documents
    are co-present. Designed to stress \emph{context bloat} (failure mode
    F1): selecting the right two documents from ten is exactly the
    \textsc{SelectChunks} problem.

  \item \textbf{Billing/Auth scenario}: a synthetic enterprise pipeline
    (5 agents: router, auth, billing, compliance, summarizer) with explicit
    PII exclusion requirements (\texttt{exclude: ["PII\_raw"]}) and
    hard \texttt{must\_carry} constraints (\texttt{user\_id},
    \texttt{order\_id}). Designed to stress \emph{silent failure}
    (failure mode F3) and the exclusion constraint (Equation~\ref{eq:excl}).
\end{itemize}

\subsection{Metrics}

Four metrics capture complementary dimensions of handoff quality:

\begin{itemize}
  \item \textbf{Must-carry recall} ($\uparrow$): fraction of
    \texttt{must\_carry} items successfully delivered to the receiving
    agent. A baseline that truncates arbitrarily will score below 1.0;
    CHP guarantees 1.0 when items are present (or raises
    \textsc{MissingMustCarry} otherwise).

  \item \textbf{Task success rate} ($\uparrow$): fraction of tasks where
    the final agent produces a correct answer against ground truth.
    Primary end-to-end quality signal.

  \item \textbf{Token efficiency} ($\downarrow$): tokens consumed at
    handoff relative to full-history baseline. Measures the cost of
    context bloat that CHP avoids.

  \item \textbf{Selection precision} ($\uparrow$): fraction of transferred
    chunks that the receiving agent actually uses in producing its output,
    oracle-judged on a 50-task subset. Measures noise in the selected set.
\end{itemize}

\subsection{Ablation Plan}

Four ablations will attribute improvements to individual protocol layers:

\begin{enumerate}[label=\textbf{A\arabic*.}]
  \item \textbf{Manifest vs.\ no manifest}: full CHP vs.\ scorer-only
    (scoring without receiver requirements) vs.\ full-history baseline.
    Isolates the contribution of Layer~1 (\texttt{ContextManifest}).

  \item \textbf{Rationale envelope impact}: task success rate with
    vs.\ without \texttt{RationaleEnvelope} injection into the receiving
    agent's context. Isolates the contribution of Layer~3.

  \item \textbf{ContextLedger recovery}: in a 5-hop chain, must-carry
    recall with full ledger access vs.\ hop-local context only. Isolates
    the contribution of Layer~4.

  \item \textbf{Scorer weight sensitivity}: grid search over
    $(\alpha, \beta, \gamma, \delta, \varepsilon)$; variance in task
    success rate across weight perturbations. Tests robustness of
    the default weight set.
\end{enumerate}

\subsection{Implementation Configuration}

The evaluation will use the following configuration:
embeddings via \texttt{sentence-transformers/all-MiniLM-L6-v2} (22M
parameters, CPU inference); token counting via \texttt{tiktoken}
(\texttt{cl100k\_base}); LLM backbone GPT-4o-mini; ContextLedger via
LanceDB 0.6 embedded. No GPU is required. All code and benchmark
configurations will be released alongside the results version of this
preprint at \url{https://github.com/[org]/chp}.

% ============================================================
\section{Discussion}
\label{sec:discussion}

\subsection{Limitations}

\textbf{Manifest quality.} The quality of CHP's selection is bounded by the
quality of the \texttt{ContextManifest} an agent publishes. Vague or
incomplete manifests produce noisy selection. We provide a manifest linting
tool and a template library for common agent types to mitigate this.

\textbf{Content-based must-carry matching.} The current implementation
matches \texttt{must\_carry} keys via substring search on chunk content. This
can fail for chunks where the key is implied rather than explicit. An
embedding-based fuzzy match is available as a configuration option.

\textbf{Scorer calibration.} Default weights are calibrated on GAIA and
HotpotQA. Different domains (code generation, scientific reasoning) may
benefit from domain-specific weight tuning.

\subsection{Future Work}

\textbf{Manifest inference.} Can a model infer an agent's likely
\texttt{ContextManifest} from its system prompt and task description? An
automatic manifest generator would lower adoption friction.

\textbf{Learned scoring.} The weighted linear scorer
(Equation~\ref{eq:score}) is interpretable but not optimal. A learned
reranker trained on \texttt{RationaleEnvelope} feedback signals could
improve precision.

\textbf{Manifest registry.} A public registry where agents advertise their
\texttt{ContextManifest} templates would enable agent discovery based on
context compatibility, not just capability.

\textbf{Formal verification.} The must-carry invariant and append-only ledger
invariant are enforced in code. Formal verification (e.g., TLA+) of the
full protocol would strengthen the safety argument.

% ============================================================
\section{Conclusion}
\label{sec:conclusion}

We have formalized agent-to-agent context handoff as a
requirements-matching problem and shown that all existing systems implement
degenerate special cases of it. The Context Handoff Protocol (CHP) introduces
the missing primitive—receiver-declared context requirements—and builds a
four-layer architecture on top of it: manifest-driven selection, rationale
provenance, and an append-only cross-hop ledger.

CHP is available as an open-source Python library and an open JSON Schema
specification, designed for drop-in adoption over existing frameworks. We
invite the multi-agent systems community to adopt, extend, and critique the
protocol. The specification and reference implementation are available at:

\begin{center}
  \url{https://github.com/[org]/chp}
\end{center}

% ============================================================
\section*{Acknowledgments}

[To be completed.]

% ============================================================
\bibliographystyle{plain}
\begin{thebibliography}{99}

\bibitem{google_a2a_2025}
Google.
\newblock Agent2Agent Protocol (A2A), 2025.
\newblock \url{https://github.com/google-a2a/A2A}.

\bibitem{anthropic_mcp_2024}
Anthropic.
\newblock Model Context Protocol, 2024.
\newblock \url{https://modelcontextprotocol.io}.

\bibitem{langgraph}
LangChain.
\newblock LangGraph: Building Stateful Multi-Agent Applications, 2024.
\newblock \url{https://github.com/langchain-ai/langgraph}.

\bibitem{crewai}
CrewAI Inc.
\newblock CrewAI: Framework for Orchestrating Role-Playing AI Agents, 2024.
\newblock \url{https://github.com/crewAIInc/crewAI}.

\bibitem{autogen}
Microsoft Research.
\newblock AutoGen: Enabling Next-Gen LLM Applications via Multi-Agent
  Conversation, 2023.
\newblock \url{https://github.com/microsoft/autogen}.

\bibitem{sk}
Microsoft.
\newblock Semantic Kernel, 2023.
\newblock \url{https://github.com/microsoft/semantic-kernel}.

\bibitem{openai_sdk}
OpenAI.
\newblock OpenAI Agents SDK, 2025.
\newblock \url{https://github.com/openai/openai-agents-python}.

\bibitem{adk}
Google.
\newblock Agent Development Kit (ADK), 2025.
\newblock \url{https://github.com/google/adk-python}.

\bibitem{packer2023memgpt}
C.~Packer et~al.
\newblock {MemGPT}: Towards LLMs as Operating Systems.
\newblock \emph{arXiv:2310.08560}, 2023.

\bibitem{mem0}
Mem0.ai.
\newblock Mem0: The Memory Layer for AI Agents, 2024.
\newblock \url{https://github.com/mem0ai/mem0}.

\bibitem{zep}
Getzep.
\newblock Zep: Long-Term Memory for AI Assistants, 2024.
\newblock \url{https://github.com/getzep/zep}.

\bibitem{baton2024}
[Author et~al.]
\newblock {BATON}: Transition-Aware Memory for Robotic Manipulation.
\newblock \emph{arXiv:2608.16889}, 2024.

\bibitem{camcp2025}
[Author et~al.]
\newblock {CA-MCP}: Collaborative Agent Extension of the Model Context
  Protocol.
\newblock \emph{arXiv:2601.11595}, 2025.

\bibitem{chap2026}
[Author et~al.]
\newblock {CHAP}: Collaborative Human-Agent Protocol.
\newblock \emph{arXiv:2606.09751}, 2026.

\bibitem{agentflow2026}
[Author et~al.]
\newblock {AgentFlow}: Static Analysis Framework for Multi-Agent Programs.
\newblock \emph{arXiv:2607.01640}, 2026.

\bibitem{reimers2019sentence}
N.~Reimers and I.~Gurevych.
\newblock Sentence-{BERT}: Sentence Embeddings using Siamese
  {BERT}-Networks.
\newblock In \emph{EMNLP}, 2019.

\bibitem{lancedb}
LanceDB.
\newblock LanceDB: Embedded Vector Database for AI Applications, 2024.
\newblock \url{https://github.com/lancedb/lancedb}.

\bibitem{mialon2023gaia}
G.~Mialon et~al.
\newblock {GAIA}: A Benchmark for General AI Assistants.
\newblock \emph{arXiv:2311.12983}, 2023.

\bibitem{yang2018hotpotqa}
Z.~Yang et~al.
\newblock {HotpotQA}: A Dataset for Diverse, Explainable Multi-hop Question
  Answering.
\newblock In \emph{EMNLP}, 2018.

\bibitem{wang2024survey}
L.~Wang et~al.
\newblock A Survey on Large Language Model based Autonomous Agents.
\newblock \emph{arXiv:2402.03578}, 2024.

\bibitem{zhang2024survey}
Z.~Zhang et~al.
\newblock A Survey on the Memory Mechanism of Large Language Model based Agents.
\newblock \emph{arXiv:2404.13501}, 2024.

\end{thebibliography}

\end{document}
